%%%%%%%%%%%%%%%%
%%% lev 16 %%%
%%%%%%%%%%%%%%%%

\Chapter{16}

\Verse{1}
{
	\hOther{וַיְדַבֵּר יְהֹוָה אֶל מֹשֶׁה אַחֲרֵי מוֹת שְׁנֵי בְּנֵי אַהֲרֹן בְּקרְבָתָם לִפְנֵי יְהֹוָה וַיָּמֻתוּ׃}
}
{
	\eOther{
		And YHWH spoke to Moses after the death of Aaron’s two
		sons—when they came forward in front of YHWH and died.
	}
}
{
%	\footnote{
%		after death. The title of this parashah is ’arê môt, meaning “after
%		death.” The title of the following week’s parashah is qdšîm, meaning
%		“holy” (19:2). In some years these two portions are read together, and
%		so their titles merge as ’arê môt qdšîm, meaning “after death: holy.” My
%		grandfather wrote a letter to my father when my father was about to be
%		married, and he noted a traditional comment on this merger: that after
%		people die we make them holy, speaking of them as if they were saints.
%		“But,” my grandfather told my father, “you have always shown me respect
%		while I am still alive.” And so one of the greatest lessons of these
%		portions of the Torah is the result of a merger of titles that are three
%		chapters apart and that were inserted centuries after the Torah was
%		completed. Thus wisdom can emerge in a variety of ways if the reader is
%		wise and sensitive.
%	}
}

\Verse{2}
{
	\hOther{וַיֹּאמֶר יְהֹוָה אֶל מֹשֶׁה דַּבֵּר אֶל אַהֲרֹן אָחִיךָ וְאַל יָבֹא בְכל עֵת אֶל הַקֹּדֶשׁ מִבֵּית לַפָּרֹכֶת אֶל פְּנֵי הַכַּפֹּרֶת אֲשֶׁר עַל הָאָרֹן וְלֹא יָמוּת כִּי בֶּעָנָן אֵרָאֶה עַל הַכַּפֹּרֶת׃}
}
{
	\eOther{
		And YHWH said to Moses, “Speak to Aaron, your brother, and
		let him not come at all times to the Holy, inside the
		pavilion, in front of the atonement dais that is over the
		ark, so he will not die, because I shall appear in a cloud
		over the atonement dais.
	}
}
{
%	\footnote{
%		inside the pavilion. The Hebrew is mibbêt lapprket. The use of the term
%		bêt (house, building) may be further evidence that the prket is a
%		four-sided structure, not a curtain. See the comments on Tabernacle
%		architecture at Exodus 26.
%	}
}

\Verse{3}
{
	\hOther{בְּזֹאת יָבֹא אַהֲרֹן אֶל הַקֹּדֶשׁ בְּפַר בֶּן בָּקָר לְחַטָּאת וְאַיִל לְעֹלָה׃}
}
{
	\eOther{
		Aaron shall come to the Holy with this: with a bull of the
		cattle for a sin offering and a ram for a burnt offering.
	}
}
{
}

\Verse{4}
{
	\hOther{כְּתֹנֶת בַּד קֹדֶשׁ יִלְבָּשׁ וּמִכְנְסֵי בַד יִהְיוּ עַל בְּשָׂרוֹ וּבְאַבְנֵט בַּד יַחְגֹּר וּבְמִצְנֶפֶת בַּד יִצְנֹף בִּגְדֵי קֹדֶשׁ הֵם וְרָחַץ בַּמַּיִם אֶת בְּשָׂרוֹ וּלְבֵשָׁם׃}
}
{
	\eOther{
		He shall wear a holy linen coat, and linen drawers shall
		be on his flesh, and he shall be belted with a linen sash,
		and he shall wear a linen headdress. They are holy
		clothes: and he shall wash his flesh with water and put
		them on.
	}
}
{
}

\Verse{5}
{
	\hOther{וּמֵאֵת עֲדַת בְּנֵי יִשְׂרָאֵל יִקַּח שְׁנֵי שְׂעִירֵי עִזִּים לְחַטָּאת וְאַיִל אֶחָד לְעֹלָה׃}
}
{
	\eOther{
		And he shall take from the congregation of the children of
		Israel two goats for a sin offering and one ram for a
		burnt offering.
	}
}
{
}

\Verse{6}
{
	\hOther{וְהִקְרִיב אַהֲרֹן אֶת פַּר הַחַטָּאת אֲשֶׁר לוֹ וְכִפֶּר בַּעֲדוֹ וּבְעַד בֵּיתוֹ׃}
}
{
	\eOther{
		And Aaron shall bring forward the bull of the sin offering
		that is his and shall make atonement for himself and for
		his house.
	}
}
{
%	\footnote{
%		his house. Meaning: his family, his household.
%	}
}

\Verse{7}
{
	\hOther{וְלָקַח אֶת שְׁנֵי הַשְּׂעִירִם וְהֶעֱמִיד אֹתָם לִפְנֵי יְהֹוָה פֶּתַח אֹהֶל מוֹעֵד׃}
}
{
	\eOther{
		And he shall take the two goats and stand them in front of
		YHWH, at the entrance of the Tent of Meeting.
	}
}
{
}

\Verse{8}
{
	\hOther{וְנָתַן אַהֲרֹן עַל שְׁנֵי הַשְּׂעִירִם גֹּרָלוֹת גּוֹרָל אֶחָד לַיהֹוָה וְגוֹרָל אֶחָד לַעֲזָאזֵל׃}
}
{
	\eOther{
		And Aaron shall cast lots on the two goats, one lot for
		YHWH and one lot for Azazel.
	}
}
{
%	\footnote{
%		lots. Why use lots? Lots are used frequently in the Tanak. In some cases
%		they are understood to reflect divine determination, and in some they
%		are not. They reflect the hand of God in the case of the lot that
%		identifies Jonah to the sailors (1:7), the lot that identifies Achan as
%		the Israelite who has violated the law of herem and taken spoils for
%		himself at Jericho (Joshua 7), the lot that identifies Saul as the first
%		king of Israel (1 Sam 10:17–21), and the lot that identifies Saul’s son
%		Jonathan as the one who has unwittingly violated an oath by his father
%		(1 Sam 14:40–42). They appear to be a matter of chance rather than
%		divine direction here in the case of the goats, in the case of the
%		apportionment of the promised land to the Israelites (Num 26:52–56;
%		Joshua 14–19; Judg 1:3), in assigning priests to their divisions (1
%		Chronicles 24), in drafting Israelites for a battle (Judg 20:9), in
%		assigning population to Jerusalem (Nehemiah 11), and in sharing in
%		spoils and captives (Joel 4:3; Obadiah 11; Nah 3:10; Ps 22:19). (The
%		case of Haman’s lots in Esther is anomalous and requires separate
%		explanation.) What is the distinction between these two groups of cases?
%		It appears to be that all the cases that involve divine determination
%		are about identifying a single individual in a specific situation. All
%		the cases that do not imply divine determination are about assigning
%		things that relate to groups or the community. The issue in the former
%		cases is finding the right person. The issue in the latter cases is
%		fairness—showing no favor in assignments. The goats are the first case
%		of such assignment by lot in the Tanak, and they are the only case that
%		is an annual event. All the others are onetime occurrences. The ritual
%		of the goats thus sets the principle and the model for all subsequent
%		cases involving lotteries for fairness in assignments. It has been said
%		that this is a denial of “Divine Providence.” But this is an act by the
%		community, seeking atonement for things that they have done. This is not
%		a situation calling for “Divine Providence.” On the contrary, the point
%		here is precisely to convey the will of the community to express their
%		atonement to God. The deity is the recipient, not the controller, of
%		this message. Footnote 16:8. Azazel. This word occurs only in this
%		chapter and nowhere else in the Tanak. There have been many attempts to
%		explain it (the common view being that it refers to some sort of desert
%		demon), but the simple fact is that we do not know who or what it is.
%	}
}

\Verse{9}
{
	\hOther{וְהִקְרִיב אַהֲרֹן אֶת הַשָּׂעִיר אֲשֶׁר עָלָה עָלָיו הַגּוֹרָל לַיהֹוָה וְעָשָׂהוּ חַטָּאת׃}
}
{
	\eOther{
		And Aaron shall bring forward the goat on which the lot
		for YHWH arose, and he shall make it a sin offering.
	}
}
{
}

\Verse{10}
{
	\hOther{וְהַשָּׂעִיר אֲשֶׁר עָלָה עָלָיו הַגּוֹרָל לַעֲזָאזֵל יעֳמַד חַי לִפְנֵי יְהֹוָה לְכַפֵּר עָלָיו לְשַׁלַּח אֹתוֹ לַעֲזָאזֵל הַמִּדְבָּרָה׃}
}
{
	\eOther{
		And the goat on which the lot for Azazel arose shall be
		stood alive in front of YHWH to make atonement over it, to
		let it go to Azazel, into the wilderness.
	}
}
{
}

\Verse{11}
{
	\hOther{וְהִקְרִיב אַהֲרֹן אֶת פַּר הַחַטָּאת אֲשֶׁר לוֹ וְכִפֶּר בַּעֲדוֹ וּבְעַד בֵּיתוֹ וְשָׁחַט אֶת פַּר הַחַטָּאת אֲשֶׁר לוֹ׃}
}
{
	\eOther{
		And Aaron shall bring forward the bull of the sin offering
		that is his and shall make atonement for him and for his
		house, and he shall slaughter the bull of the sin offering
		that he has.
	}
}
{
}

\Verse{12}
{
	\hOther{וְלָקַח מְלֹא הַמַּחְתָּה גַּחֲלֵי אֵשׁ מֵעַל הַמִּזְבֵּחַ מִלִּפְנֵי יְהֹוָה וּמְלֹא חפְנָיו קְטֹרֶת סַמִּים דַּקָּה וְהֵבִיא מִבֵּית לַפָּרֹכֶת׃}
}
{
	\eOther{
		And he shall take a fire-holder-full of coals of fire from
		on the altar, from in front of YHWH, and handfuls of fine
		incense of fragrances, and he shall bring it inside the
		pavilion
	}
}
{
}

\Verse{13}
{
	\hOther{וְנָתַן אֶת הַקְּטֹרֶת עַל הָאֵשׁ לִפְנֵי יְהֹוָה וְכִסָּה עֲנַן הַקְּטֹרֶת אֶת הַכַּפֹּרֶת אֲשֶׁר עַל הָעֵדוּת וְלֹא יָמוּת׃}
}
{
	\eOther{
		and put the incense on the fire in front of YHWH so that
		the cloud of the incense will cover the atonement dais
		that is over the Testimony, and he will not die.
	}
}
{
}

\Verse{14}
{
	\hOther{וְלָקַח מִדַּם הַפָּר וְהִזָּה בְאֶצְבָּעוֹ עַל פְּנֵי הַכַּפֹּרֶת קֵדְמָה וְלִפְנֵי הַכַּפֹּרֶת יַזֶּה שֶׁבַע פְּעָמִים מִן הַדָּם בְּאֶצְבָּעוֹ׃}
}
{
	\eOther{
		And he shall take some of the bull’s blood and sprinkle
		with his finger on the front of the atonement dais
		eastward, and he shall sprinkle some of the blood with his
		finger in front of the atonement dais seven times.
	}
}
{
}

\Verse{15}
{
	\hOther{וְשָׁחַט אֶת שְׂעִיר הַחַטָּאת אֲשֶׁר לָעָם וְהֵבִיא אֶת דָּמוֹ אֶל מִבֵּית לַפָּרֹכֶת וְעָשָׂה אֶת דָּמוֹ כַּאֲשֶׁר עָשָׂה לְדַם הַפָּר וְהִזָּה אֹתוֹ עַל הַכַּפֹּרֶת וְלִפְנֵי הַכַּפֹּרֶת׃}
}
{
	\eOther{
		And he shall slaughter the goat of the sin offering that
		is the people’s and shall bring its blood inside the
		pavilion and shall do its blood as he did to the bull’s
		blood: and he shall sprinkle it on the atonement dais and
		in front of the atonement dais.
	}
}
{
}

\Verse{16}
{
	\hOther{וְכִפֶּר עַל הַקֹּדֶשׁ מִטֻּמְאֹת בְּנֵי יִשְׂרָאֵל וּמִפִּשְׁעֵיהֶם לְכל חַטֹּאתָם וְכֵן יַעֲשֶׂה לְאֹהֶל מוֹעֵד הַשֹּׁכֵן אִתָּם בְּתוֹךְ טֻמְאֹתָם׃}
}
{
	\eOther{
		And he shall make atonement over the Holy from the
		impurities of the children of Israel and from their
		offenses, for all their sins, and he shall do so for the
		Tent of Meeting that tents with them among their
		impurities.
	}
}
{
}

\Verse{17}
{
	\hOther{וְכל אָדָם לֹא יִהְיֶה בְּאֹהֶל מוֹעֵד בְּבֹאוֹ לְכַפֵּר בַּקֹּדֶשׁ עַד צֵאתוֹ וְכִפֶּר בַּעֲדוֹ וּבְעַד בֵּיתוֹ וּבְעַד כּל קְהַל יִשְׂרָאֵל׃}
}
{
	\eOther{
		And no human shall be in the Tent of Meeting when he comes
		to make atonement in the Holy until he comes out. And he
		shall make atonement for him and for his house and for all
		the community of Israel.
	}
}
{
}

\Verse{18}
{
	\hOther{וְיָצָא אֶל הַמִּזְבֵּחַ אֲשֶׁר לִפְנֵי יְהֹוָה וְכִפֶּר עָלָיו וְלָקַח מִדַּם הַפָּר וּמִדַּם הַשָּׂעִיר וְנָתַן עַל קַרְנוֹת הַמִּזְבֵּחַ סָבִיב׃}
}
{
	\eOther{
		“And he shall go out to the altar that is in front of YHWH
		and make atonement over it, and he shall take some of the
		bull’s blood and some of the goat’s blood and put it on
		the altar’s horns all around.
	}
}
{
}

\Verse{19}
{
	\hOther{וְהִזָּה עָלָיו מִן הַדָּם בְּאֶצְבָּעוֹ שֶׁבַע פְּעָמִים וְטִהֲרוֹ וְקִדְּשׁוֹ מִטֻּמְאֹת בְּנֵי יִשְׂרָאֵל׃}
}
{
	\eOther{
		And he shall sprinkle some of the blood on it with his
		finger seven times and purify it and make it holy from the
		impurities of the children of Israel.
	}
}
{
}

\Verse{20}
{
	\hOther{וְכִלָּה מִכַּפֵּר אֶת הַקֹּדֶשׁ וְאֶת אֹהֶל מוֹעֵד וְאֶת הַמִּזְבֵּחַ וְהִקְרִיב אֶת הַשָּׂעִיר הֶחָי׃}
}
{
	\eOther{
		“And when he will finish making atonement for the Holy and
		the Tent of Meeting and the altar, then he shall bring
		forward the live goat.
	}
}
{
}

\Verse{21}
{
	\hOther{וְסָמַךְ אַהֲרֹן אֶת שְׁתֵּי יָדָו עַל רֹאשׁ הַשָּׂעִיר הַחַי וְהִתְוַדָּה עָלָיו אֶת כּל עֲוֺנֹת בְּנֵי יִשְׂרָאֵל וְאֶת כּל פִּשְׁעֵיהֶם לְכל חַטֹּאתָם וְנָתַן אֹתָם עַל רֹאשׁ הַשָּׂעִיר וְשִׁלַּח בְּיַד אִישׁ עִתִּי הַמִּדְבָּרָה׃}
}
{
	\eOther{
		And Aaron shall lay his two hands on the live goat’s head
		and confess over it all the crimes of the children of
		Israel and all their offenses, for all their sins, and he
		shall put them on the goat’s head and let it go by the
		hand of an appointed man to the wilderness.
	}
}
{
}

\Verse{22}
{
	\hOther{וְנָשָׂא הַשָּׂעִיר עָלָיו אֶת כּל עֲוֺנֹתָם אֶל אֶרֶץ גְּזֵרָה וְשִׁלַּח אֶת הַשָּׂעִיר בַּמִּדְבָּר׃}
}
{
	\eOther{
		And the goat will carry all their crimes on it to an
		inaccessible land. And he shall let the goat go in the
		wilderness.
	}
}
{
}

\Verse{23}
{
	\hOther{וּבָא אַהֲרֹן אֶל אֹהֶל מוֹעֵד וּפָשַׁט אֶת בִּגְדֵי הַבָּד אֲשֶׁר לָבַשׁ בְּבֹאוֹ אֶל הַקֹּדֶשׁ וְהִנִּיחָם שָׁם׃}
}
{
	\eOther{
		“And Aaron shall come to the Tent of Meeting and take off
		the linen clothes that he wore when he came to the Holy,
		and he shall leave them there.
	}
}
{
}

\Verse{24}
{
	\hOther{וְרָחַץ אֶת בְּשָׂרוֹ בַמַּיִם בְּמָקוֹם קָדוֹשׁ וְלָבַשׁ אֶת בְּגָדָיו וְיָצָא וְעָשָׂה אֶת עֹלָתוֹ וְאֶת עֹלַת הָעָם וְכִפֶּר בַּעֲדוֹ וּבְעַד הָעָם׃}
}
{
	\eOther{
		And he shall wash his flesh with water in a holy place and
		put on his clothes, and he shall go out and make his burnt
		offering and the people’s burnt offering and make
		atonement for him and for the people.
	}
}
{
}

\Verse{25}
{
	\hOther{וְאֵת חֵלֶב הַחַטָּאת יַקְטִיר הַמִּזְבֵּחָה׃}
}
{
	\eOther{
		And he shall burn the fat of the sin offering into smoke
		at the altar.
	}
}
{
}

\Verse{26}
{
	\hOther{וְהַמְשַׁלֵּחַ אֶת הַשָּׂעִיר לַעֲזָאזֵל יְכַבֵּס בְּגָדָיו וְרָחַץ אֶת בְּשָׂרוֹ בַּמָּיִם וְאַחֲרֵי כֵן יָבוֹא אֶל הַמַּחֲנֶה׃}
}
{
	\eOther{
		And the one who let the goat go to Azazel shall wash his
		clothes and wash his flesh with water, and after that he
		shall come to the camp.
	}
}
{
}

\Verse{27}
{
	\hOther{וְאֵת פַּר הַחַטָּאת וְאֵת שְׂעִיר הַחַטָּאת אֲשֶׁר הוּבָא אֶת דָּמָם לְכַפֵּר בַּקֹּדֶשׁ יוֹצִיא אֶל מִחוּץ לַמַּחֲנֶה וְשָׂרְפוּ בָאֵשׁ אֶת עֹרֹתָם וְאֶת בְּשָׂרָם וְאֶת פִּרְשָׁם׃}
}
{
	\eOther{
		And he shall take the bull of the sin offering and the
		goat of the sin offering, whose blood was brought to make
		atonement in the Holy, outside the camp; and they shall
		burn their skin and their meat and their dung in fire.
	}
}
{
}

\Verse{28}
{
	\hOther{וְהַשֹּׂרֵף אֹתָם יְכַבֵּס בְּגָדָיו וְרָחַץ אֶת בְּשָׂרוֹ בַּמָּיִם וְאַחֲרֵי כֵן יָבוֹא אֶל הַמַּחֲנֶה׃}
}
{
	\eOther{
		And the one who burns them shall wash his clothes and wash
		his flesh with water, and after that he shall come to the
		camp.
	}
}
{
}

\Verse{29}
{
	\hOther{וְהָיְתָה לָכֶם לְחֻקַּת עוֹלָם בַּחֹדֶשׁ הַשְּׁבִיעִי בֶּעָשׂוֹר לַחֹדֶשׁ תְּעַנּוּ אֶת נַפְשֹׁתֵיכֶם וְכל מְלָאכָה לֹא תַעֲשׂוּ הָאֶזְרָח וְהַגֵּר הַגָּר בְּתוֹכְכֶם׃}
}
{
	\eOther{
		“And it shall be an eternal law for you: in the seventh
		month, on the tenth of the month, you shall degrade
		yourselves, and you shall not do any work: the citizen and
		the alien who resides among you.
	}
}
{
%	\footnote{
%		degrade yourselves. The verb here is the same word that is used to
%		describe what the Egyptians did to the Israelites (Exod 1:11–12):
%		degradation, affliction. It has been understood to mean that one does
%		not eat or drink or bathe or wear fine shoes on the Day of Atonement;
%		rather, one humbles oneself.
%	}
}

\Verse{30}
{
	\hOther{כִּי בַיּוֹם הַזֶּה יְכַפֵּר עֲלֵיכֶם לְטַהֵר אֶתְכֶם מִכֹּל חַטֹּאתֵיכֶם לִפְנֵי יְהֹוָה תִּטְהָרוּ׃}
}
{
	\eOther{
		Because on this day he shall make atonement over you, to
		purify you. You will be pure from all your sins in front
		of YHWH.
	}
}
{
}

\Verse{31}
{
	\hOther{שַׁבַּת שַׁבָּתוֹן הִיא לָכֶם וְעִנִּיתֶם אֶת נַפְשֹׁתֵיכֶם חֻקַּת עוֹלָם׃}
}
{
	\eOther{
		It is a Sabbath, a ceasing, to you. And you shall degrade
		yourselves—an eternal law.
	}
}
{
}

\Verse{32}
{
	\hOther{וְכִפֶּר הַכֹּהֵן אֲשֶׁר יִמְשַׁח אֹתוֹ וַאֲשֶׁר יְמַלֵּא אֶת יָדוֹ לְכַהֵן תַּחַת אָבִיו וְלָבַשׁ אֶת בִּגְדֵי הַבָּד בִּגְדֵי הַקֹּדֶשׁ׃}
}
{
	\eOther{
		And the priest whom one will anoint and who will fill his
		hand to function as a priest in his father’s place shall
		make atonement. And he shall wear the linen clothes, the
		holy clothes.
	}
}
{
%	\footnote{
%		fill his hand to function as a priest. See comment on Exod 29:29.
%	}
}

\Verse{33}
{
	\hOther{וְכִפֶּר אֶת מִקְדַּשׁ הַקֹּדֶשׁ וְאֶת אֹהֶל מוֹעֵד וְאֶת הַמִּזְבֵּחַ יְכַפֵּר וְעַל הַכֹּהֲנִים וְעַל כּל עַם הַקָּהָל יְכַפֵּר׃}
}
{
	\eOther{
		And he shall make atonement for the holy place of the
		Holy, and he shall make atonement for the Tent of Meeting
		and for the altar, and he shall make atonement over the
		priests and over all the people of the community.
	}
}
{
}

\Verse{34}
{
	\hOther{וְהָיְתָה זֹּאת לָכֶם לְחֻקַּת עוֹלָם לְכַפֵּר עַל בְּנֵי יִשְׂרָאֵל מִכּל חַטֹּאתָם אַחַת בַּשָּׁנָה וַיַּעַשׂ כַּאֲשֶׁר צִוָּה יְהֹוָה אֶת מֹשֶׁה׃}
}
{
	\eOther{
		And this shall become an eternal law for you, to make
		atonement over the children of Israel from all their sins,
		once in the year.” And he did as YHWH had commanded Moses.
	}
}
{
}
